% Осень26!B36; нумерация внутри столбца без пустых строк.
% Источник: лекция 7, §7 (строка 218)
% Источник: лекция 7, §8 (строка 235)
\definitionticket{33}{Сходимость по мере}
\label{definition:33}

\subsection*{Сходимость по мере}

\begin{definition}

\emph{сходимость по мере.} Пусть \(f_n,f\colon X\to\overline{\mathbb R}\) --- измеримые, \textbf{почти везде конечные} функции (множества бесконечных значений имеют меру ноль). Говорят, что \textbf{\(f_n\) сходится к \(f\) по мере \(\mu\)} (обозначение лектора: «толстая стрелочка» с пометкой \(\mu\) снизу; здесь: \(f_n\xrightarrow[\mu]{}f\)), если \[\forall\varepsilon>0:\qquad \mu\bigl(\bigl\{x\in X:\ |f_n(x)-f(x)|\ge\varepsilon\bigr\}\bigr)\ \longrightarrow\ 0\qquad(n\to\infty).\] «На макроязыке»: мера множества, где \(f_n\) сильно отличается от \(f\) (не меньше чем на \(\varepsilon\)), стремится к нулю.

\end{definition}

\textbf{Связь с теорией вероятностей.} В вероятностных терминах это \textbf{сходимость по вероятности} (пространство элементарных событий с вероятностью --- мера, \(\mu(X)=1\)). Отличия нашего курса: (i) не задаём вопросов теории вероятностей; (ii) наши меры бывают неограниченными --- «как раз наоборот, сплошь и рядом \(\lambda(\mathbb R^m)=+\infty\)».

\textbf{Замечание 1 (нечувствительность к замене предела на множестве меры нуля).} Если поменять \(f\) на множестве \(N\) меры ноль, то при каждом \(n,\varepsilon\) множество \(\{|f_n-f|\ge\varepsilon\}\) изменится лишь на часть \(N\) меры ноль --- мера не изменится, сходимость по мере сохраняется. Значит, предел по мере определён \textbf{однозначно лишь с точностью до эквивалентности} \(f\sim g\).

\textbf{Философское отступление лектора (функции как фактор-классы).} Функция, которую «имеешь право заменить на множестве меры ноль», --- «странный объект»: значением в отдельной точке пользоваться нельзя (в любой точке имеешь право его заменить). «Это полностью выключает мозг: что же мы тогда можем и в каком смысле функция дана?» Ответ: воспринимать надо \textbf{не функцию, а фактор-класс по отношению эквивалентности} \(\sim\). В точках класс «значения не имеет», но \textbf{интегрировать его можно}: взять любого представителя класса и проинтегрировать; замена представителя меняет подынтегральную функцию лишь на множестве меры ноль, интеграл «этого не почувствует». К этому «надо привыкнуть... чуть попозже мы ещё раз на этом намедитируем».

\textbf{Замечание 2 (единственность предела по мере).} Если \(f_n\xrightarrow[\mu]{}f\) и \(f_n\xrightarrow[\mu]{}g\), то \(f\sim g\) (\(f=g\) п.в.). \textbf{Требует сверки.} Сформулировано лектором без доказательства («тогда то отличие, которое я заметил, только оно и может быть»). Стандартное обосновнение (не звучало): \(|f-g|\le|f-f_n|+|f_n-g|\), поэтому \(\{|f-g|\ge2\varepsilon\}\subset\{|f-f_n|\ge\varepsilon\}\cup\{|f_n-g|\ge\varepsilon\}\), и по полуаддитивности мера левого множества нулевая при всех \(\varepsilon\); останется заметить \(\{f\ne g\}=\bigcup_k\{|f-g|\ge 1/k\}\).

\subsection*{Три примера о соотношении сходимостей}

Речь о соотношении: (а) поточечная сходимость (всюду), (б) сходимость п.в., (в) сходимость по мере. Примеры показывают, что сходимости «находятся в общем положении».

\textbf{Пример 1 (обе сходимости есть).} \(X=(0,+\infty)\) (можно и \(\mathbb R\)), мера Лебега \(\lambda\). Пусть \[f_n(x)=\frac{1}{n\,x}.\] Динамика прозрачна: \(f_1\) --- школьная гипербола \(1/x\), \(f_2=\tfrac12\cdot\tfrac1x\) («два раза прижали»), \(f_3=\tfrac13\cdot\tfrac1x\) и т.д. --- графики всё льнут к оси \(x\). Тогда:

\begin{itemize}
\tightlist
\item
  \textbf{поточечно} \(f_n(x)\to0\) всюду на \((0,+\infty)\) (на \(\mathbb R\) --- всюду, кроме нуля);
\item
  \textbf{по мере} тоже \(f_n\to0\): «в лоб» по определению, \[\{x:\ f_n(x)\ge\varepsilon\}=\Bigl(0,\ \tfrac{1}{n\varepsilon}\Bigr],\qquad \lambda\bigl(\{f_n\ge\varepsilon\}\bigr)=\frac{1}{n\varepsilon}\xrightarrow[n\to\infty]{}0\qquad\forall\varepsilon>0.\] Итог: есть и сходимость п.в., и сходимость по мере --- обе к нулю.
\end{itemize}

\textbf{Пример 2 (есть сходимость п.в., нет сходимости по мере).} \(X=\mathbb R\), \(\lambda\). Пусть \[f_n(x)=e^{-(x-n)^2}\] --- «колокольчик» \(e^{-x^2}\), целиком переезжающий в точку \(n\) (форма при сдвиге не меняется). Тогда:

\begin{itemize}
\tightlist
\item
  \textbf{поточечно} \(f_n(x)\to0\) при каждом фиксированном \(x\): когда \(n\) станет больше \(x\), колокольчик уедет, а по формуле показатель \(-(x-n)^2\to-\infty\);
\item
  \textbf{по мере к нулю сходимости нет}: посмотрим, где \(f_n\ge\tfrac12\) (проведём горизонталь на высоте \(1/2\)): это интервал \textbf{одной и той же ширины} вокруг точки \(n\), т.е. \[\lambda\bigl(\{f_n\ge\tfrac12\}\bigr)=\mathrm{const}>0\qquad\text{не зависит от }n\] --- пик фиксированной ширины просто едет вправо. По определению сходимости по мере эта мера должна была бы стремиться к нулю. Сходимости к нулю нет. \textbf{Требует сверки.} К функциям, отличным от нуля, сходимость по мере лектор не анализировал: «кажется, кроме нуля больше некуда сходить» (неформальная эвристика: на любом фиксированном отрезке при больших \(n\) функция \(f_n\) почти неотличима от нуля).
\end{itemize}

\textbf{Мораль.} На пространстве \textbf{бесконечной меры} сходимость почти везде \textbf{не влечёт} сходимость по мере: «масса» уезжает на бесконечность, не убывая.

\textbf{Пример 3 («блуждающие ступеньки»: есть сходимость по мере, нет сходимости даже поточечно).} \(X=[0,1]\), \(\lambda\). Каждое \(n\in\mathbb N\) запишем в виде \[n=2^{k}+l,\qquad k\ge0,\ \ 0\le l<2^{k}\] (в двоичной записи --- «отщепляем первую единицу»: \(2^k\) --- старший разряд, остаток --- слово \(l\)). Положим \[f_n:=\chi_{\left[\,l/2^{k},\ (l+1)/2^{k}\right)}.\] \emph{Устройство («счётчик»):} при фиксированном \(k\) ступенька ширины \(2^{-k}\) перебирает положения \(l=0,1,\dots,2^{k}-1\), пробегая слева направо весь \([0,1]\); затем \(k\) увеличивается на единицу, отрезочки становятся вдвое короче (их вдвое больше) --- и счётчик идёт заново.

\begin{itemize}
\tightlist
\item
  \textbf{Поточечного предела нет ни в одной точке \(x\in[0,1]\):} «вы сидите в точке \(x\) и ждёте»: при каждом очередном \(k\) ступенька добежит до вас --- будет \(1\), --- и уйдёт --- снова \(0\); значения \(0\) и \(1\) принимаются бесконечно много раз.
\item
  \textbf{По мере \(f_n\to0\):} для любого \(0<\varepsilon\le1\) \[\lambda\bigl(\{f_n\ge\varepsilon\}\bigr)=\lambda\Bigl(\bigl[\,l/2^{k},(l+1)/2^{k}\bigr)\Bigr)=\frac{1}{2^{k}}\xrightarrow[n\to\infty]{}0\] (при \(n\to\infty\) имеем \(k\to\infty\)); при \(\varepsilon>1\) множество пусто.
\end{itemize}

\textbf{Мораль.} Сходимость по мере \textbf{не влечёт} сходимость почти везде. Процесс «затухает в целом»: места, где он «топорщится», локализованы на маленьком блуждающем отрезочке, и с вероятностной точки зрения это неважно («подарок от теории вероятностей: такое бывает»). Итог по трём примерам: сходимости п.в. и по мере «в общем положении» --- бывают обе, бывает одна без другой и другая без первой. «Сюрприз.»
